% Methods: water classification and the active zone (A3). Design: D-015, D-027, D-029 (zone v6), D-048; PREREG section 3 and Amendments 3 and 9.
% Citation still to add once the key exists: Sentinel-2 L2A product documentation for the scene classification (SCL).

\subsection*{Water classification}
\emph{Active zone.} Scenes were classified only within an active zone, a band meant to contain the waterlines of the altimetry era. Its bounds are the water extents of the \metric{params.highstand_label} high stand and the \metric{params.lowstand_label} low stand, mapped from clear May--September observations. For these two extents only, a pixel counts as water if its MNDWI (below) is above zero, or if its near-infrared and first shortwave-infrared reflectances are both below \metric{params.dark_water_reflectance}: clear, deep water in the south of the domain is dark in both bands, and its MNDWI scatters around zero. The landward bound is the near-maximum extent of the high stand: pixels that were water in at least \metric{params.high_extent_min_water_pct}\% of clear Landsat~5 observations between \metric{params.highstand_start} and \metric{params.highstand_end}. Where fewer than \metric{params.min_obs_high} such observations exist, the maximum extent of the Global Surface Water record\citep{pekel2016} takes their place. The seaward bound is the near-minimum extent of the low stand: pixels that were water in at least \metric{params.low_extent_min_water_pct}\% of clear Landsat~8--9 observations between \metric{params.lowstand_start} and \metric{params.lowstand_end}. Only the part of the band between the two extents that is connected to the open sea is kept. The open sea is low-stand water whose surroundings within \metric{params.sea_radius_km}~km are at least \metric{params.sea_min_water_pct}\% water, which excludes lakes. High-stand water is kept if a path through high-stand water, at most \metric{params.conn_max_km}~km long, joins it to the open sea. The paths are traced on a \metric{params.conn_scale_km}~km grid on which a cell counts as water if any of its pixels does, so that narrow channels stay connected. The connected band was dilated by \metric{params.zone_dilate_m}~m on a \metric{params.zone_scale_m}~m grid in the equal-area projection. Per-pixel computation was limited to the \metric{params.zone_tile_km}~km cells that contain the zone. The low-stand bound reflects the levels of the May--September scenes; lower levels were reached in other months, so whether the zone contains the waterlines of the lowest level bins is checked for each level and sector before any hindcast comparison.

\emph{Snow and ice.} Two tests flag snow or ice in cloud-free pixels. The first is the product quality layer: the snow bit of the Landsat quality band\citep{usgs_l45tm_c2l2,usgs_l7etm_c2l2,usgs_l89oli_c2l2}, or Sentinel-2 scene-classification class \metric{params.s2_scl_snow}. The second is the conservative snow test of the Theia snow algorithm\citep{gascoin2019}: a normalised difference snow index above \metric{params.ndsi_snow} together with red reflectance above \metric{params.red_snow}. The red-reflectance condition keeps turbid water from being flagged as snow. Pixels flagged by either test were removed from \metric{params.snow_season_first_month} to \metric{params.snow_season_last_month}. A test was applied in every month only if, in each May--September month, it flagged at most \metric{params.snow_rule_max_pct}\% of the clear pixels in the zone (pre-registration Amendments~3 and~9). The share flagged by each test is recorded for every scene and month.

\emph{Threshold.} For each scene we computed the modified normalised difference water index (MNDWI) from the green and first shortwave-infrared bands\citep{xu2006}. Its water threshold was set by Otsu's method\citep{otsu1979}, as the value that maximises the between-class variance of the scene's MNDWI histogram. The histogram has \metric{params.otsu_bins} equal bins between \metric{params.mndwi_min} and \metric{params.mndwi_max} and is sampled at \metric{params.hist_scale_m}~m inside the active zone. Sampling inside the zone replaces the edge-guided sample of the pre-registered design, at far lower computing cost (pre-registration Amendment~3). The edge-guided threshold is computed for a random subset of scenes as a robustness check. Each threshold is computed once and stored, so the per-pixel classification is a single comparison.

\emph{Classes.} Pixels above the threshold are water. Pixels below it with a normalised difference vegetation index above \metric{params.ndvi_veg} are emergent vegetation; these reed beds hide the water's edge, so they are treated as unobserved. All remaining pixels are dry. Overlapping scenes of one overpass (adjacent Landsat rows, or Sentinel-2 tiles of one datatake) were mosaicked first, so that each overpass counts at most once per pixel. The automated water extraction index AWEI\textsubscript{sh}\citep{feyisa2014} and a fixed MNDWI threshold are the alternatives tested in the robustness checks.
